\section{Proposed approach: GARDIAN}
\label{sec:method}

\subsection{Overview of GARDIAN}
\label{sec:method-overview}

GARDIAN is a lightweight re-ranking layer for hybrid retrieval in medical
Retrieval-Augmented Generation (RAG). Placed between a hybrid retriever and a
frozen reader LLM, it reorders the candidate passages returned by sparse and
dense retrieval before they are handed to the reader. Two properties
distinguish it from conventional fusion. First, both retrieval channels are
scored by \emph{learned} branch heads over within-pool--normalised retrieval
statistics, rather than being combined at their raw score scales. Second, the
weight given to each channel is predicted \emph{per query} by a controller,
rather than fixed globally as in Sum or RRF. GARDIAN never re-encodes
query--passage pairs, so it costs a small fraction of a cross-encoder pass.

\begin{figure}[t]
  \centering
  \includegraphics[width=\linewidth]{fig_training.pdf}
  \caption{\textbf{Offline training.} The retrieval back-ends are frozen; only
  the two branch heads and the controller are trained. Each query contributes
  one candidate \emph{pool} rather than isolated pairs, and the listwise
  objective weights every comparison by the nDCG@$k$ change of swapping it, so
  gradient concentrates on pairs that straddle the cutoff.}
  \label{fig:training}
\end{figure}

\begin{figure}[t]
  \centering
  \includegraphics[width=\linewidth]{fig_inference.pdf}
  \caption{\textbf{Online inference.} The controller runs \emph{once} per
  query on the query embedding alone; the branch heads score cached
  first-stage features. No query--passage re-encoding occurs at any point,
  which is what separates GARDIAN's latency from that of cross-encoder
  re-rankers.}
  \label{fig:inference}
\end{figure}

Figures~\ref{fig:training} and~\ref{fig:inference} show the offline and online
paths. A query is first processed independently by a sparse and a dense
retriever, whose top-$50$ lists are merged into a unified candidate pool. For
each query--passage pair GARDIAN extracts lexical and semantic features that
are scored by two dedicated branches. In parallel, a controller predicts the
relative importance of the two channels from the query representation alone.
The predicted weights combine the branch scores into a final relevance score,
and the top-$K$ re-ranked passages are forwarded to the reader.

\paragraph{Task formulation.}
Given a query $q$, the hybrid retriever constructs a candidate pool
$\mathcal{P}(q)=\{p_1,\ldots,p_{K'}\}$ from the union of the sparse and dense
result lists, capped at $K'\!=\!100$. Each candidate carries a score from each
channel, or a missing-score marker where that channel did not retrieve it.
GARDIAN learns a scoring function assigning a relevance score $r_i$ to every
candidate, so that $\mathcal{P}(q)$ can be re-ranked and its top-$K$ passages
passed to the frozen reader.

\subsection{Branch scoring}
\label{sec:features}

GARDIAN scores lexical and semantic evidence through two separate branches
rather than concatenating all signals into one representation. Keeping the
channels apart is what makes a per-channel fusion weight meaningful: the
controller can re-weight one channel without disturbing the other.

Both feature vectors live in $\mathbb{R}^{8}$. The \textbf{sparse branch}
receives the BM25 or SPLADE$^{++}$ score, IDF-weighted term overlap, Jaccard
overlap, the within-pool min--max and $z$-normalised score, the within-pool
rank, the channel's score dispersion, and an indicator of whether this channel
retrieved the candidate at all. The \textbf{dense branch} receives the
corresponding dense quantities: the FAISS or MedCPT score, mean and maximum
absolute embedding difference, within-pool $z$ and min--max normalisations,
within-pool rank, dispersion, and its own retrieval indicator.

The normalisation features matter more than the raw scores. A union pool takes
the top-$k$ of each channel, so a substantial fraction of candidates carry no
score from one of them; recording that as $0$ would assert a falsehood, since
the smallest genuinely measured similarity is well above zero. All
normalisations are therefore computed over \emph{scored} candidates only, with
unscored ones imputed at the channel floor and flagged by the indicator. In
our feature ablation (Section~\ref{sec:ablation}) the four normalisation
features carry the majority of the learned model's attributed gain, while the
raw retrieval scores carry almost none.

Each branch applies a three-layer MLP~\cite{goodfellow2016deep} with LayerNorm
to its feature vector, producing a scalar estimate:
%
\begin{equation}
  \hat{s}_i^{\mathrm{sp}} = f_{\mathrm{sp}}\!\left(\mathbf{x}_i^{\mathrm{sp}}\right),
  \qquad
  \hat{s}_i^{\mathrm{de}} = f_{\mathrm{de}}\!\left(\mathbf{x}_i^{\mathrm{de}}\right),
  \qquad
  \mathbf{x}_i^{\mathrm{sp}},\, \mathbf{x}_i^{\mathrm{de}} \in \mathbb{R}^{8}.
  \label{eq:branches}
\end{equation}
%
The two branches share this architecture but not their parameters. The
estimates remain separate until the controller combines them, preserving the
complementarity of lexical and semantic evidence~\cite{burges2005ranknet}.

\subsection{Query-adaptive controller}
\label{sec:controller}

The controller predicts how much each retrieval channel should contribute,
from the query alone and before any candidate is examined. It takes a single
input: the mean-pooled contextual embedding
$\mathbf{h}_q \in \mathbb{R}^{768}$ produced by
PubMedBERT~\cite{gu2021pubmedbert}. This is passed through a three-layer MLP
whose two outputs are normalised by a softmax:
%
\begin{equation}
  (\alpha_s,\, \alpha_d)
  \;=\;
  \mathrm{softmax}\!\big(\mathrm{MLP}(\mathbf{h}_q)\big),
  \qquad \alpha_s + \alpha_d = 1.
  \label{eq:controller}
\end{equation}
%
The softmax imposes a \emph{simplex constraint}: the model allocates a fixed
budget between the two channels rather than amplifying both, which prevents
the degenerate solution in which the distinction between channels collapses
into an overall score scaling.

The controller is conditioned on the query embedding only. An earlier version
also received a one-hot question-type indicator; we removed it after measuring
that the indicator is constant within each benchmark and therefore separates
the datasets rather than the queries, making it a dataset identifier rather
than a query signal. Question type is retained purely as an \emph{analysis}
dimension and is never a model input.

Because the controller runs once per query and its weights are shared across
all $K'$ candidates, adaptive fusion adds one encoder pass per query and
leaves re-ranking complexity at $\mathcal{O}(K')$. Section~\ref{sec:ablation}
quantifies how much this adaptivity is worth against a fixed weight tuned on
development data; we report that comparison rather than assuming the benefit.

\subsection{Score fusion}
\label{sec:fusion}

Given Equations~(\ref{eq:branches}) and~(\ref{eq:controller}), the final
relevance score of candidate $p_i$ is the convex combination
%
\begin{equation}
  r_i
  \;=\;
  \alpha_s\,\hat{s}_i^{\mathrm{sp}}
  \;+\;
  \alpha_d\,\hat{s}_i^{\mathrm{de}}.
  \label{eq:fusion}
\end{equation}
%
Fixing $\alpha_s=\alpha_d=\tfrac{1}{2}$ recovers uniform averaging over the
learned branches, and any fixed $\alpha_s$ recovers weighted hybrid
fusion~\cite{cormack2009rrf} over them. GARDIAN generalises both by predicting
$(\alpha_s,\alpha_d)$ per query. Note that Equation~(\ref{eq:fusion}) combines
\emph{learned} branch estimates, not raw retrieval scores: this distinction is
what separates GARDIAN from a tuned $\alpha$-fusion baseline, and we isolate
the two effects in Section~\ref{sec:ablation}.

\subsection{Offline training}
\label{sec:offline}

Only the two branch heads and the controller are trained. The sparse and dense
retrievers, the PubMedBERT query encoder, and the reader LLM stay frozen.
Query embeddings $\mathbf{h}_q$ are computed once with mean pooling and cached.
One model is trained per hybrid family, independently, on five seeds
$\{13,21,42,87,100\}$. Training concatenates the PubMedQA-Artificial and
MedMCQA \emph{train} splits of that family's rank JSONL; PubMedQA-Labeled has
no train split and is evaluation-only. Model selection uses nDCG@10 on the
concatenated development splits of the same two collections.

\paragraph{Training pool.}
Rank JSONL stores the first-stage union pool in retrieval order. A query
contributes a training example only if it has at least one labelled positive
and one negative. All positives are kept (at most $63$), and the pool is
filled to $N{=}64$ candidates. Of those negatives, $85\%$ are drawn from the
first $120$ negatives in first-stage order and the rest from the remainder, so
the loss sees top-$k$ confusions rather than easy tail passages. The $N$ rows
are then shuffled: pool index must not leak the first-stage rank through
padding position. Each step scores $B{=}64$ such queries. A reservoir shuffle
of $20{,}000$ emitted pools mixes the two collections, which otherwise appear
as one PubMedQA-Artificial block followed by one MedMCQA block.

\paragraph{Forward pass.}
For a query $q$ with candidates $i=1,\ldots,N$, the frozen encoder supplies
$\mathbf{h}_q\in\mathbb{R}^{768}$. Each candidate has an $8$-dimensional sparse
vector and an $8$-dimensional dense vector (Section~\ref{sec:features}). The
two BranchMLPs map $8\to 2048\to 1024\to 1$ with LayerNorm, ReLU and dropout
$0.12$ after each hidden layer, and do not share weights. The controller is
the same shape with a two-dimensional output, takes $\mathbf{h}_q$ only, and
applies a softmax, giving $(\alpha_s,\alpha_d)$ once per query. The score is
Equation~(\ref{eq:fusion}). Branch outputs are \emph{not} re-standardised
before mixing.

\paragraph{Objective.}
The configured loss is LambdaRank~\cite{burges2010ranknet} at cutoff
$k{=}10$. Let $y_i\in\{0,1\}$ be the stored grade, $G_i=2^{y_i}-1$ (equal to
$y_i$ on binary labels), and $r_i$ the current score. Ranks are obtained by
sorting $r$ with padding masked off, and are treated as constants
(no gradient through the sort). With $0$-based rank $\rho_i$,
$D_i=1/\log_2(\rho_i+2)$ and $\mathrm{IDCG}@k$ computed from the ideal
ordering of $G$ truncated at $k$, the pairwise weight is
%
\begin{equation}
  w_{ij}
  \;=\;
  \mathbf{1}[y_i > y_j]\;
  \frac{|G_i-G_j|\;|D_i-D_j|}{\mathrm{IDCG}@k}.
  \label{eq:lambdarank-w}
\end{equation}
%
The batch loss is the mean, over queries that have a positive and at least two
live candidates, of the weighted logistic
%
\begin{equation}
  \mathcal{L}(\theta)
  \;=\;
  \frac{1}{|Q_{\mathrm{valid}}|}
  \sum_{q\in Q_{\mathrm{valid}}}
  \sum_{i,j}
  w_{ij}^{q}\;
  \mathrm{softplus}\!\big(-\sigma\,(r_i^{q}-r_j^{q})\big),
  \qquad \sigma = 1.
  \label{eq:loss}
\end{equation}
%
Queries with no positive contribute nothing and are excluded from the
denominator. This is the optimiser's objective: a per-pair average is logged
but is not what AdamW minimises. A pairwise margin
$\mathrm{softplus}(\gamma-(r^+-r^-))$ is implemented in the same trainer and
is not used.

\paragraph{Optimisation.}
AdamW jointly updates both branches and the controller
($\mathrm{lr}=10^{-3}$, weight decay $1.5\times10^{-4}$). The learning rate
warms up linearly over the first three epochs, then follows a cosine decay to
$0.08\times\mathrm{lr}$. Gradients are clipped to global norm $1.0$. The
forward pass uses mixed precision; Equation~(\ref{eq:loss}) is evaluated in
full precision. Training runs for at most $10$ epochs and stops if development
nDCG@10 does not improve for $6$ epochs. The saved checkpoint is the epoch
with the best development nDCG@10.

\begin{figure}[t]
\centering
\fbox{\begin{minipage}{0.97\linewidth}\small
\textbf{Algorithm 1.} GARDIAN offline training for one hybrid family and one
seed.\\[0.4em]
\textbf{Input:} rank JSONL (frozen retrievers), cached $\mathbf{h}_q$, seed
$s$.\\
\textbf{Output:} checkpoint with best development nDCG@10.
\begin{enumerate}\setlength{\itemsep}{2pt}
\item Concatenate PubMedQA-Artificial train and MedMCQA train; concatenate
      their development splits for selection.
\item For each query with $\geq 1$ positive and $\geq 1$ negative: keep all
      positives (cap $63$); fill to $N{=}64$ with negatives ($85\%$ from the
      first $120$ first-stage negatives, $15\%$ from the rest); shuffle the
      $N$ rows.
\item Mini-batch: $B{=}64$ queries. Sparse/dense BranchMLP $\to$
      $\hat s^{\mathrm{sp}},\hat s^{\mathrm{de}}$. Controller
      $\mathrm{softmax}(\mathrm{MLP}(\mathbf{h}_q))\to(\alpha_s,\alpha_d)$.
      $r_i=\alpha_s\hat s_i^{\mathrm{sp}}+\alpha_d\hat s_i^{\mathrm{de}}$.
\item Compute $w_{ij}$ from the current ranking with no gradient
      (Eq.~\ref{eq:lambdarank-w}); minimise Eq.~\ref{eq:loss} with AdamW,
      grad-norm clip $1.0$.
\item After each epoch, score the development pools with the same forward
      pass. Keep the best nDCG@10; stop after $6$ non-improving epochs, or
      after $10$ epochs.
\end{enumerate}
\end{minipage}}
\caption{Training procedure as implemented. Retrieval, the query encoder and
the reader are never updated. Five seeds and four hybrid families are trained
independently with this algorithm.}
\label{fig:alg-train}
\end{figure}

\subsection{Online inference}
\label{sec:online}

At inference GARDIAN is a thin layer between retrieval and generation. Given
$\mathcal{P}(q)$, the controller predicts the fusion weights once per query
from Equation~(\ref{eq:controller}); the branch heads score every candidate
from cached first-stage features; Equation~(\ref{eq:fusion}) combines them;
and the top-$K$ passages go to the frozen reader.

Two properties keep this cheap. Re-ranking cost is linear in $|\mathcal{P}(q)|$
and involves only two small MLP passes per candidate over pre-computed
features. And no query--passage pair is ever re-encoded, unlike a
cross-encoder which must run a full transformer forward pass for each of the
$K'$ candidates. The only neural encoding GARDIAN performs at query time is
the single embedding $\mathbf{h}_q$ required by the controller.
